\documentclass[10pt,a4paper]{amsart}
\usepackage[margin=19mm]{geometry}
\usepackage{amsmath,amssymb,mathtools}
\usepackage[scr=rsfs,cal=euler]{mathalfa}
\usepackage{xfrac}
\usepackage[unicode,hidelinks,bookmarks=false]{hyperref}
\newcommand{\ie}{i.e.\ }
\newcommand{\cf}{cf.\ }
\newcommand{\resp}{respectively\ }
\newcommand{\Subsection}{\subsection}
\providecommand{\nobreakdash}{\nobreak\hbox{-}}
\setlength{\emergencystretch}{2em}
\multlinegap0pt
\usepackage{amssymb}
\usepackage{xcolor}
\usepackage{tikz}
\newtheorem{theorem}{Theorem}
\newtheorem{lemma}{Lemma}
\newtheorem{proposition}{Proposition}
\title{Equilibrium configurations of a\\ 3D fluid-beam interaction problem}
\author{Vincenzo Bianca, Edoardo Bocchi, Filippo Gazzola}
\date{Source: arXiv:2507.10523v3 (CC BY 4.0)}
\begin{document}
\maketitle

\noindent\emph{Source and licence.} Equilibrium configurations of a\\ 3D fluid-beam interaction problem. Version arXiv:2507.10523v3 (CC BY 4.0), distributed by arXiv under the Creative Commons Attribution 4.0 International licence, http://creativecommons.org/licenses/by/4.0/, as recorded in the arXiv licence metadata for this exact version and shown on the abstract page https://arxiv.org/abs/2507.10523v3. Full licence notice: \texttt{licenses/arxiv-2507.10523-CC-BY-4.0.txt}.\par\smallskip

\noindent\textit{Editorial scope.} Counted unit: the article's proposition L-incre (source0.tex lines 776--781). The article's complete original proof of it is delivered with it (source0.tex lines 782--912, one \texttt{proof} environment of 131 lines). No mathematics is written, repaired or re-derived: the statement and the proof are the article's own lines, in the article's own notation, and no retained line is substituted or re-flowed.  Retained uncounted as the source-local context the proof needs: lines 264--266; lines 326--328; lines 398--406; lines 408--413; lines 427--431; lines 428--430; lines 609--611; lines 641--647; lines 720--732; lines 760--764; lines 765--772.  Nothing was omitted from any retained span, so the package carries no omission record.  No retained line contains a citation, so the wrapper carries no bibliography and no bibliography entry.  No retained line contains a figure environment, an image-inclusion command or drawing code, so no image is delivered and the package declares no asset dependency.  Editorial formatting only, in addition: an AMS \texttt{amsart} preamble that restates the article's own declarations and macros used by the retained lines, namely the environments \texttt{theorem}, \texttt{lemma}, \texttt{proposition} on separate counters, together with the standard packages those lines need.  The retained environments print in this document as Theorem 1, Lemma 1, Proposition 1 in order of appearance, whereas the article prints the same items with the numbering of its own full document.  The complete native source of the article as distributed in the arXiv e-print for this exact version is bundled unchanged as \texttt{sources/181572/source0.tex}.
\begin{equation}\label{eq_bound_infty}
\left\|h\right\|_{L^\infty(-1,1)} <	1- \omega\max\{Z^+,Z^- \}.\\[5pt]
\end{equation}

	\begin{equation}\label{lift}
		L(y, u,p)= -e_3\cdot\int_{\partial\mathcal{B}_y^h} \mathbb{T}(u, p)n= -e_3\cdot\int_{\partial\mathcal{B}_y^h}[\eta(\nabla u+\nabla^Tu)-p\mathbb{I}]n,
	\end{equation}

\begin{equation}\label{u*}
u_*=
\begin{cases}
0&\hspace{.3in}\mbox{on}\hspace{.1in}\partial\mathcal{B}^h\\
\gamma Ve_1&\hspace{.3in}\mbox{on}\hspace{.1in}\overline{\Gamma_{y}^\pm \setminus \mathcal{B}_{\pm 1}}\cup\Gamma_z^\pm\\
\gamma V_ie_1&\hspace{.3in}\mbox{on}\hspace{.1in}\Gamma_x^-\\
\gamma V_oe_1&\hspace{.3in}\mbox{on}\hspace{.1in}\Gamma_x^+.
\end{cases}
\end{equation}

\begin{equation}\label{eq_main}
\begin{aligned}
&-\eta\Delta u+u\cdot\nabla u+\nabla p=0,&\qquad&\nabla\cdot u=0&\qquad&\mbox{in}\hspace{.1in}\Omega_h\\
&u=u_*&\qquad& &\qquad&\mbox{on}\hspace{.1in}\partial\Omega_h.
\end{aligned}
\end{equation}

\begin{theorem}\label{the_ex_regu}Let $h\in C_0^{1,1}([-1,1])$ be small enough to ensure \eqref{eq_bound_infty}. For any $\gamma\geq0$, \eqref{eq_main} admits a weak solution $(u, p)$ in $\Omega_h$; there exists $\overline{\gamma}>0$ such that the weak solution is unique for $\gamma\in [0, \overline{\gamma})$. Moreover, any weak solution is a strong solution and there exists $C^*>0$, depending on $\sigma$ but not on $h$, such that
\begin{equation}\label{reg-est}
\left\|u\right\|_{W^{2,2+\sigma }(\Omega_h)}+\left\|p\right\|_{W^{1,2+\sigma}(\Omega_h)}\le C^*\gamma \qquad \forall  \gamma\in[0, \overline{\gamma}).
\end{equation}
\end{theorem}

\begin{equation}
\left\|u\right\|_{W^{2,2+\sigma }(\Omega_h)}+\left\|p\right\|_{W^{1,2+\sigma}(\Omega_h)}\le C^*\gamma \qquad \forall  \gamma\in[0, \overline{\gamma}).
\end{equation}

    \begin{equation}\label{load}
        L(y, h)=L(y, u(h), p(h)) = - e_3 \cdot \int_{\partial \mathcal{B}_y^h}\mathbb{T}(u(h), p(h)) n
    \end{equation} with $\mathbb{T}(\cdot, \cdot)$ as in \eqref{lift}.

\begin{lemma}\label{lem-Lbound}
Let $h\in C^{1,1}_0([-1,1])$ satisfy the smallness assumption in Theorem \ref{the_ex_regu}. Let $\gamma \in [0, \overline{\gamma})$ and $(u,p)\in W^{2,2+\sigma}(\Omega_h)\times W^{1,2+\sigma}(\Omega_h)$ be the unique solution to \eqref{eq_main} from Theorem \ref{the_ex_regu}. Then, there exists  $C>0$, independent of $h$,  such that
\begin{equation}\label{load-control}
\|L(\cdot, h)\|_{L^\infty(-1,1)}\leq C \gamma .
\end{equation}
Moreover, $L(\cdot, h)\in C(-1,1)$.
\end{lemma}

\begin{lemma}\label{lemma-diffeo}
Let $\mathcal{D}\subset \Gamma_y^-$ be an open set such that the cylinder $\mathcal{C}=\mathcal{D} \times [-1,1]$ strictly contains any $\mathcal{B}^h$ for $h$ sufficiently small. Consider a smooth cut-off function $\xi:[-R, R]\times [-1,1]\rightarrow [0,1]$ such that
\begin{equation}\label{cutoff-diffeo}\begin{aligned}\xi=1\quad  \text{in} \quad \overline{\mathcal{D}} \quad \text{and} \quad
\xi=0 \quad \text{in}\quad \left([-R\times R]\times [-1,1]\right)\setminus \lambda\mathcal{D}
    \end{aligned}
\end{equation}with $\lambda>1$ ensuring that the homothetic $\lambda \mathcal{C}\subset \Gamma_y^-$.
For $h\in \mathbb{H}$, consider the map
	\begin{equation}\label{diffeo}\varphi_h (x,y,z)= \left(x, y, z+ \xi(x,z)h(y)\right).\end{equation} Then, there exists $\beta>0$ such that, for $\|h\|_{\mathbb{H}}<\beta$,  $\varphi_h$ is a $C^1$-diffeomorphism from $\overline{\Omega}_0$ to $\overline{\Omega}_h$ satisfying \begin{equation}\label{diff-prop}
			\varphi_h (\partial\mathcal{B})=\partial\mathcal{B}_h , \quad \varphi_h (\overline{\Gamma_x^\pm })=\overline{\Gamma_x^\pm},
           \quad \varphi_h (\overline{\Gamma_y^\pm \setminus \mathcal{B}_{\pm 1} })=\overline{\Gamma_y^\pm \setminus \mathcal{B}_{\pm 1} } \ \
              \text{and} \ \ \varphi_h (\overline{\Gamma_z^\pm })=\overline{\Gamma_z^\pm}.
            \end{equation}
\end{lemma}

\begin{equation}\label{matrices-diffeo}
\begin{aligned}
    A_h= \det(J_h) M_h^TM_h \quad \text{and} \quad B_h= \det(J_h) M_h \quad \text{with} \quad  M_h = (J_h^{-1})^T,\\[5pt]
\end{aligned}
\end{equation}we obtain that $(U,P)=(u, p)\circ \varphi_h$ solves the elliptic problem

\begin{equation}\label{BVprob_trans}
		\begin{aligned}
			&-\eta\nabla \cdot\left( A_h \nabla U \right)  +  U\cdot B_h\nabla U+ B_h\nabla P=0, \quad
			\nabla\cdot (B_h^TU)=0\quad\text{in}\quad \Omega_0\\[2pt]
			&U=0 \quad \text{on} \quad  \partial \mathcal{B} ,\quad U=\gamma Ve_1\quad \text{on} \quad \overline{\Gamma_{y}^\pm \setminus \mathcal{B}_{\pm 1}}\cup\Gamma_z^\pm\\[2pt]
			&U= \gamma V_ie_1\quad \text{on} \quad \Gamma_x^-,\quad U= \gamma V_oe_1\quad\text{on} \quad \Gamma_x^+.\\[5pt]
		\end{aligned}
	\end{equation}

\begin{proposition}\label{L-incre}
	There exists $\beta_0\in (0, \beta]$ with $\beta$ as in Lemma \ref{lemma-diffeo} such that $ h\mapsto L( y, h)$ is Lipschitz continuous in $\mathbb{B}_{\beta_0}$. More precisely, for any $h_1,h_2 \in \mathbb{B}_{\beta_0}$,
\begin{equation}\label{lip-lift}
    \|L(\cdot, h_1) - L(\cdot, h_2)\|_{L^\infty(-1,1)} \leq  \gamma C(\beta_0) \|h_1-h_2\|_{\mathbb{H}^4}
\end{equation}with $\gamma \in [0,\gamma_0)$ for some $\gamma_0\in(0,\overline{\gamma}]$ and $\overline{\gamma}$ as in Theorem \ref{the_ex_regu}.
\end{proposition}

\begin{proof}
With $\varphi_h$ at hand,   the lift in \eqref{load}
can be written in terms of the $h$-independent fluid domain $\Omega_0$, namely,
    \begin{equation} \label{lift-trans}\begin{aligned}
        L(y,h)&=-e_3\cdot \int_{\partial \mathcal{B}_y}\widetilde{\mathbb{T}}_h(U, P) M_hn_0\\&=-e_3 \cdot \int_{\partial \mathcal{B}_y}  \left[\eta (M_h \nabla U + (M_h\nabla U)^T) - P\mathbb{I}\right] M_hn_0
        \end{aligned}
    \end{equation}
where $(U,P)$ solves \eqref{BVprob_trans} and $n_0$ is the  unit normal vector to $\partial \mathcal{B}$ pointing inside $\mathcal{B}$. In \eqref{lift-trans}, we have used the fact $\det(J_h)=1$ on $\partial \mathcal{B}_y$. We now show that, there exists $\beta_0>0$ such that
\begin{equation}\label{claim}\vspace{0.1em}
h\mapsto (U(h),P(h))\in W^{2, 2+\sigma}(\Omega_0) \times W^{1, 2+\sigma}(\Omega_0)\mbox{ belongs to }  C^1(\mathbb{B}_{\beta_0}).\vspace{0.1em}
\end{equation}
 To this end, with $\beta$ as in Lemma \ref{lemma-diffeo} we express \eqref{BVprob_trans} as
\begin{equation}\label{implicit}
\mathcal{H}(h,U, P)=0,
\end{equation}where  $\mathcal{H} :  \mathbb{B}_\beta\times  X \rightarrow Y$  with
\begin{equation*}\begin{aligned}
&X=W^{2, 2+\sigma}(\Omega_0)\times W^{1, 2+\sigma}(\Omega_0),\\
&Y=\bigg\{(\chi_1,\chi_2,\chi_3)\in L^{2+\sigma}(\Omega_0)\times W^{1, 2+\sigma}(\Omega_0)\times W^{2-\tfrac{1}{2+\sigma}, 2+\sigma}(\partial \Omega_0)  \\
&\qquad \qquad  \text{such that} \quad \int_{\Omega_0} \chi_2 = \int_{\partial \Omega_0}B_h^T \chi_3 \cdot n \bigg\}, \end{aligned}\end{equation*} and the components of $\mathcal{H}$ read
\begin{equation}\label{defH}
	\begin{aligned}
	\mathcal{H}_1(h, U, P)= &-\eta\nabla \cdot (A_h\nabla U) + U \cdot B_h\nabla U + B_h \nabla P,\\[5pt]
	\mathcal{H}_2(h, U, P)= & \ \nabla \cdot (B_h^T U), \\[5pt]
    \mathcal{H}_3(h, U, P)= & \ U_{|_{\partial \Omega_0}} - u_*,
	\end{aligned}
\end{equation}where $u_*$ is as in \eqref{u*}. We explicitly compute the matrices in \eqref{matrices-diffeo}, which read
\begin{equation}\label{exp-matrices}\begin{aligned}
 &A_h= \mathbb{I}  +\left(\begin{matrix}
\partial_z \xi h
  & 0 & -\partial_x \xi h \\[5pt]0&\partial_z \xi h & - \xi h'\\[5pt]
-\partial_x \xi h& - \xi h'&\frac{-\partial_z\xi h + (\partial_x \xi h)^2 + (\xi h')^2}{1+ \partial_z \xi h}
\end{matrix}\right), \
 &B_h= \mathbb{I} +\left(\begin{matrix}
\partial_z \xi h  & 0&-\partial_x \xi h \\[5pt]0 & \partial_z \xi h&  - \xi h'\\[5pt]
0 & 0& 0
\end{matrix}\right).\\[5pt]
\end{aligned}
\end{equation}
First, we have that $h\mapsto\mathcal{H}(h, U, P)$ is $C^1(\mathbb{B}_{\beta})$ since Lemma \ref{lemma-diffeo} yields the same regularity for the maps $h\mapsto A_h$, $h\mapsto B_h$ and $h\mapsto B_h^T$. Second, $(U,P)\mapsto \mathcal{H}(h, U,P)$ is $C^1(X)$ for any $h\in \mathbb{B}_{\beta}$; indeed, the linear differential operator $\mathcal{L}=D_{(U, P)}\mathcal{H}[h,U,P]$, whose components read
\begin{equation*}\begin{aligned}
\mathcal{L}_1( \chi, \Pi)= &-\eta\nabla \cdot (A_h\nabla \chi) + \chi \cdot B_h\nabla U + U \cdot B_h\nabla \chi+ B_h\nabla \Pi,\\[5pt]
\mathcal{L}_2( \chi, \Pi)= & \ \nabla \cdot (B_h^T\chi),\\[5pt]
\mathcal{L}_3( \chi, \Pi)  = &  \ \chi_{|_{\partial \Omega_0}},
\end{aligned}
\end{equation*} is bounded from  $X$ to $Y$, hence it is Fréchet differentiable. Moreover, it is continuous with respect to $(h,U,P)\in \mathbb{B}_\beta\times X$.
Thus, $(h, U, P) \mapsto\mathcal{H}(h, U, P)$ belongs to $C^1(  \mathbb{B}_\beta\times X)$. \par Moreover, we claim that $\mathcal{L}$ is an isomorphism from $X$ to $Y$. To prove this assertion, let us consider $(\phi_1, \phi_2, \phi_3 )\in Y$ and the linear elliptic problem \begin{equation}\label{iso-eq}\mathcal{L}(\chi, \Pi)= (\phi_1, \phi_2, \phi_3), \end{equation} which can be written as the inhomogeneous Stokes-type system
\begin{equation}\label{stokes-mod}\begin{aligned}
&-\eta \Delta \chi + \nabla \Pi =  - \chi \cdot B_h\nabla U -U\cdot B_h\nabla \chi \\
&  \quad \qquad \qquad \qquad-\eta\nabla \cdot ((\mathbb{I}-A_h)\nabla \chi) + (\mathbb{I}-B_h)\nabla \Pi + \phi_1 \quad &&&&\mbox{in}  \quad \Omega_0 \\[5pt]\ &\nabla\cdot \chi  = \ \nabla \cdot (\mathbb{I}-B_h^T \chi) +\phi_2 \quad &&&&\mbox{in} \quad \Omega_0 \\[5pt]
&\chi_{|_{\partial \Omega_0}}= \ \phi_3.
\end{aligned}
\end{equation}
By \eqref{exp-matrices} and since $h\in \mathbb{B}_\beta$, there exist $C_1(\beta),C_2(\beta)>0$, with $C_{1}(\beta), C_2(\beta)\rightarrow 0$ as $\beta\rightarrow 0$, such that
$$\|\mathbb{I}-A_h\|_{W^{1, \infty}(\Omega_0)}\leq C_1(\beta) \quad \text{and} \quad \|\mathbb{I}-B_h\|_{W^{2,\infty}(\Omega_0)}\leq C_2(\beta).$$ We hence infer that
there exist $\beta_0\in (0, \beta]$ and $r_0>0$ such that, by a contraction argument,  \eqref{stokes-mod} admits a unique solution $(\chi, \Pi)\in X$ provided that
\begin{equation}\label{small-cond}
\|h\|_{\mathbb{H}^4}<\beta_0 \qquad \mbox{and} \quad 	\|U\|_{W^{2, 2+\sigma}(\Omega_0)}< r_0.
\end{equation} Moreover, the properties of the diffeomorpshim $\varphi_h$ imply that
\begin{equation}\label{controlU-u}\begin{aligned}
 c\|u\|_{W^{2, 2+\sigma}(\Omega_h)}\leq &\|U\|_{W^{2, 2+\sigma}(\Omega_0)}\leq C\|u\|_{W^{2, 2+\sigma}(\Omega_h)},\\[5pt] c\|p\|_{W^{1, 2+\sigma}(\Omega_h)}\leq &\|P\|_{W^{1, 2+\sigma}(\Omega_0)}\leq C\|p\|_{W^{1, 2+\sigma}(\Omega_h)},
 \end{aligned}
\end{equation} with constants $0<c\leq C$ independent of $h$, by taking $\gamma \in (0,\overline{\gamma})$ sufficiently small, the bound \eqref{reg-est}
yields the second condition in \eqref{small-cond}. Thus, $\mathcal{L}$ is an isomorphism and there exists $C>0$, independent of $h$, such that the unique solution to \eqref{iso-eq} satisfies
\begin{equation}\label{est-iso}\begin{aligned}
    &\|\chi\|_{W^{2, 2+\sigma}(\Omega_0)} + \|\Pi\|_{W^{1, 2+\sigma}(\Omega_0)} \\[5pt]&\leq C \left( \|\phi_1\|_{L^{2+\sigma}(\Omega_0)} + \|\phi_2\|_{W^{1, 2+\sigma}(\Omega_0)} + \|\phi_3\|_{W^{2-\frac{1}{2+\sigma}, 2+\sigma}(\partial \Omega_0)}\right)\end{aligned}
\end{equation}
Therefore,  applying the Implicit Function Theorem to \eqref{implicit}, we conclude \eqref{claim}.

In addition, it follows from \eqref{claim} that there differential operators $D_h U[h] : \mathbb{B}_{\beta_0}\rightarrow W^{2, 2+\sigma}(\Omega_0)$ and $D_h P[h] : \mathbb{B}_{\beta_0}\rightarrow W^{1,2+\sigma}(\Omega_0)$ are bounded. We also know from \eqref{implicit} that, for any $\tilde h\in \mathbb{B}_{\beta_0}$,  $$U_h(\tilde{h}):=D_h U[h] \tilde{h}\quad \text{and} \quad P_h(\tilde{h}):=D_h P[h] \tilde{h}$$ solve the linear elliptic problem
\begin{equation*}
\mathcal{L}( U_h(\tilde{h}), P_h(\tilde{h}) )= - D_h\mathcal{H}[h, U, P]\tilde{h},
\end{equation*} which explicitly reads
\begin{equation*}
\begin{aligned}
&-\eta\nabla \cdot (A_h\nabla U_h(\tilde{h})) + U_h(\tilde{h}) \cdot B_h\nabla U + U \cdot B_h\nabla U_h(\tilde{h}) + B_h\nabla P_h(\tilde{h})= F(h, \tilde{h}, U, P)
&\mbox{in} \ \Omega_0 \\[5pt]
&\nabla \cdot (B_h^T U_h(\tilde{h})) = G(h, \tilde{h}, U) &\mbox{in} \ \Omega_0\\[5pt]
&U_h(\tilde{h})_{|_{\partial \Omega_0}}=0
\end{aligned}
\end{equation*}
with the source terms
\begin{equation*}\begin{aligned}
    &F(h, \tilde{h}, U, P)=-\eta \nabla \cdot (A_h(\tilde{h})\nabla U) + U\cdot B_h(\tilde{h})\nabla U + B_h(\tilde{h})\nabla P,\\[5pt]&
    G(h, \tilde{h}, U)= \nabla \cdot ( B^T_h(\tilde{h}) U)
    \end{aligned}
\end{equation*}
and
\begin{equation*}
A_h(\tilde{h})=- D_h A_h[h]\tilde{h}, \quad
B_h(\tilde{h})= -D_h B_h[h]\tilde{h},\quad
B_h^T(\tilde{h})= -D_h B_h^T[h]\tilde{h}, \\[5pt]
\end{equation*} where $D_h A_h[h]$, $ D_h B_h[h]$ and $D_h B_h^T[h]$ are the differential operators associated with the maps $h\mapsto A_h$, $h\mapsto B_h$ and $h\mapsto B_h^T$.
Since $\mathcal{L}$ is an isomorphism, we know that $(U_h(\tilde{h}), P_h(\tilde{h}))$ is uniquely determined and, by \eqref{est-iso}, it satisfies
\begin{equation}\label{est-appdiff}\begin{aligned}
    \|U_h(\tilde{h})\|_{W^{2, 2+\sigma}(\Omega_0)} &+ \|P_h(\tilde{h})\|_{W^{1, 2+\sigma}(\Omega_0)} \\[5pt]&\leq C \left( \|F(h, \tilde{h}, U, P)\|_{L^{2+\sigma}(\Omega_0)} + \|G(h, \tilde{h}, U)\|_{W^{1, 2+\sigma}(\Omega_0)} \right).\end{aligned}
\end{equation}
For any $h\in \mathbb{B}_{\beta_0}$, it follows from the explicit expressions \eqref{exp-matrices} that there exists $C(\beta_0)>0$ such that
\begin{equation*}\begin{aligned}
    \|A_h(\tilde{h})\|_{H^2(\Omega_0)}\leq C(\beta_0)&\|\tilde{h}\|_{\mathbb{H}^4}, \qquad   \|B_h(\tilde{h})\|_{H^2(\Omega_0)}\leq C(\beta_0)\|\tilde{h}\|_{\mathbb{H}^4},
    \end{aligned}
\end{equation*} so that, using the continuous embedding $H^2(\Omega_0)\subset L^\infty(\Omega_0)$, the right-hand side of \eqref{est-appdiff} can be estimated by
\begin{equation}\label{est-sources-dif}
    \begin{aligned}
        &\|F(h, \tilde{h}, U, P)\|_{L^{2+\sigma} (\Omega_0)} + \|G(h, \tilde{h}, U)\|_{W^{1,2+\sigma} (\Omega_0)}\\[5pt]&\leq C\left(\|A_h(\tilde{h})\|_{H^2(\Omega_0)}+\|B_h(\tilde{h})\|_{H^2(\Omega_0)}\right)\left( \| U\|_{W^{2, 2+\sigma}(\Omega_0)} + \| P\|_{W^{1, 2+\sigma}(\Omega_0) } \right)
     \\[5pt]&  \leq  C(\beta_0)\|\tilde{h}\|_{\mathbb{H}^4}\left(\| U\|_{W^{2, 2+\sigma}(\Omega_0)} + \| P\|_{W^{1, 2+\sigma}(\Omega_0)} \right).
    \end{aligned}
\end{equation}
Due to \eqref{controlU-u} and \eqref{reg-est}, we then obtain
\begin{equation*}
    \begin{aligned}
    \|U_h(\tilde{h})\|_{W^{2, 2+\sigma}(\Omega_0)} + \|P_h(\tilde{h})\|_{W^{1, 2+\sigma}(\Omega_0)} \leq C(\beta_0)\gamma  \|\tilde{h}\|_{\mathbb{H}^4},
    \end{aligned}
\end{equation*}which, in turn, implies the uniform bound for the norms of the differential operators $D_hU[h]$ and $D_hP[h]$
\begin{equation}\label{uni-bou-diff}
    \sup_{h\in \mathbb{B}_{\beta_0}} \|D_hU[h]\| +  \sup_{h\in \mathbb{B}_{\beta_0}} \|D_hP[h]\| \leq C(\beta_0)\gamma.
\end{equation}

Finally, let us estimate the variation of the lift, defined in \eqref{lift-trans}, with respect to the variation of $h$. Using the same argument and the Trace Theorem as in the proof of Lemma \ref{lem-Lbound}, we have that,  for any $h_1,h_2\in \mathbb{B}_{\beta_0}$,
\begin{equation}\label{diff-lift}\begin{aligned}
&|	L(y, h_1) - 	L(y, h_2) |\\[5pt]&=
\bigg|\int_{\partial \mathcal{B}_y}  \left(\widetilde{\mathbb{T}}_{h_1}(U(h_1), P(h_1)) M_{h_1}n_0 - \widetilde{\mathbb{T}}_{h_2}(U(h_2), P(h_2)) M_{h_2}n_0\right) \bigg|
\\[5pt]&
=\bigg| \int_{\partial \mathcal{B}_y}\bigg[\left(\widetilde{\mathbb{T}}_{h_1}(U(h_1),P(h_1)) - \widetilde{\mathbb{T}}_{h_2}(U(h_2),P(h_2)) \right) M_{h_1}n_0 \\[2pt]&\qquad \qquad\qquad  +\widetilde{\mathbb{T}}_{h_2}(U(h_2),P(h_2)) (M_{h_1} - M_{h_2})n_0\bigg]\bigg|\\[5pt]&
\leq C(\beta_0)\bigg(\|U(h_1)- U(h_2)\|_{W^{2, 2+\sigma}(\Omega_0)}+ \|P(h_1)-P(h_2)\|_{W^{1, 2+\sigma}(\Omega_0)} \\[2pt]&\qquad\qquad  +\left(\|U(h_2)\|_{W^{2,2+\sigma}(\Omega_0)}+ \|P(h_2)\|_{W^{1, 2+\sigma}(\Omega_0)}\right)\|h_1-h_2\|_{\mathbb{H}^4}\bigg).\\[5pt]
\end{aligned}
\end{equation}
 Combining the Mean Value Inequality with \eqref{uni-bou-diff} and using \eqref{controlU-u} together with \eqref{reg-est}, we obtain that
\begin{equation*}\begin{aligned} &\|	L(\cdot, h_1) - 	L(\cdot, h_2) \|_{L^\infty(-1,1)}\\[5pt]&\leq C(\beta_0)\left(\gamma + \|u(h_2)\|_{W^{2, 2+\sigma}(\Omega_{h_2})}+ \|p(h_2)\|_{W^{1, 2+\sigma}(\Omega_{h_2})}\right)\|h_1-h_2\|_{\mathbb{H}^4}\\[5pt]&\leq \gamma C(\beta_0)\|h_1-h_2\|_{\mathbb{H}^4},
\end{aligned}
\end{equation*} which proves the Lipschitz continuity of $L(\cdot, h)$ with respect to $h$.
\end{proof}

\end{document}
